\documentclass[11pt]{article}
\usepackage[margin=1in]{geometry}
\usepackage{amsmath,amssymb,amsthm}
\usepackage{graphicx}
\usepackage[hidelinks]{hyperref}

\newtheorem{theorem}{Theorem}
\newtheorem{lemma}{Lemma}
\newtheorem{corollary}{Corollary}

\newcommand{\N}{\mathbb{N}}
\newcommand{\Per}{\operatorname{Per}}
\newcommand{\ran}{\operatorname{ran}}
\newcommand{\dist}{\operatorname{dist}}

\title{A Closed-Range Subcase of Dense Small Periodic Sets}
\author{Run \texttt{fa\_banach\_001}, lane 4}
\date{2026-06-29}

\begin{document}
\maketitle

\section*{Source And Target Question}

The source paper is Rodrigo Cardeccia and Santiago Muro,
\emph{Arithmetic progressions and chaos in linear dynamics},
arXiv:2003.07161.  In Question 4.2 the authors ask:

\begin{quote}
Is any hypercyclic operator with dense small periodic sets necessarily chaotic?
Or, more generally, does any operator with dense small periodic sets have dense
periodic points?
\end{quote}

\begin{center}
\includegraphics[width=0.92\linewidth]{figures/open_problem_crop.png}
\end{center}

The packet proves a positive Banach-space subcase under a closed-range
hypothesis on the periodic equations.

\section*{Definitions}

Let $T$ be a bounded linear operator on a Banach space $X$. A set
$Y\subset X$ is a periodic set for $T$ if $T^k(Y)\subset Y$ for some
$k\geq 1$. The operator $T$ has dense small periodic sets if every nonempty
open set $U\subset X$ contains a closed periodic set.

We use the equivalent form from Proposition 3.4 of Cardeccia--Muro: for every
nonempty open set $U$ there are $k\geq 1$ and $x\in U$ such that
$T^{jk}x\in U$ for every $j\geq 1$.

\section*{Proof Intuition}

The source paper turns small periodic sets into actual periodic points by
placing the periodic set inside a compact convex hull and applying a fixed
point theorem. The argument below replaces compactness by a closed-range
metric estimate.

If an orbit of $S=T^k$ stays forever in a small ball, then the Cesaro averages
of that orbit stay in the same ball and satisfy
\[
  (I-S)z_N = \frac{x-S^N x}{N}.
\]
In a Banach space this tends to zero because the orbit is bounded. Thus each
small ball contains approximate fixed points of $S$. When $I-S$ has closed
range, approximate fixed points of $S$ must lie close to genuine fixed points
of $S$. Those fixed points are periodic points for $T$.

\begin{lemma}[Closed-range estimate]
Let $A:X\to X$ be a bounded linear operator on a Banach space. If $\ran A$ is
closed, then there is a constant $C>0$ such that
\[
  \dist(y,\ker A)\leq C\|Ay\| \qquad (y\in X).
\]
\end{lemma}

\begin{proof}
The operator induced by $A$,
\[
  \widetilde A:X/\ker A\longrightarrow \ran A,\qquad
  \widetilde A(y+\ker A)=Ay,
\]
is a bounded bijection between Banach spaces. Since $\ran A$ is closed, the
open mapping theorem gives a bounded inverse. Hence
\[
  \|y+\ker A\|_{X/\ker A}\leq C\|Ay\|.
\]
This quotient norm is exactly $\dist(y,\ker A)$.
\end{proof}

\begin{theorem}
Let $T\in\mathcal L(X)$ be a bounded linear operator on a Banach space $X$.
Assume that $T$ has dense small periodic sets and that $\ran(I-T^n)$ is closed
for every $n\geq 1$. Then $\Per(T)$ is dense in $X$.
\end{theorem}

\begin{proof}
Let $U\subset X$ be nonempty and open. Choose $x_0\in X$ and $r>0$ such that
$B(x_0,r)\subset U$. Apply dense small periodicity to the smaller ball
$B(x_0,r/2)$. There are $k\geq 1$ and $x\in B(x_0,r/2)$ such that, writing
$S=T^k$, one has
\[
  S^j x\in B(x_0,r/2)\qquad (j=0,1,2,\ldots).
\]

For $N\geq 1$ set
\[
  z_N=\frac1N\sum_{j=0}^{N-1}S^j x .
\]
Because the ball $B(x_0,r/2)$ is convex, each $z_N$ belongs to
$B(x_0,r/2)$. Moreover
\[
  (I-S)z_N=\frac{x-S^N x}{N},
\]
and therefore
\[
  \|(I-S)z_N\|
  \leq \frac{\|x-x_0\|+\|S^N x-x_0\|}{N}
  <\frac{r}{N}.
\]

By hypothesis, $I-S=I-T^k$ has closed range. The closed-range estimate gives
a constant $C>0$ such that
\[
  \dist(z_N,\ker(I-S))\leq C\|(I-S)z_N\|<\frac{Cr}{N}.
\]
Choose $N$ so large that $Cr/N<r/2$. Then some $p\in\ker(I-S)$ satisfies
$\|p-z_N\|<r/2$. Since $z_N\in B(x_0,r/2)$, we get $p\in B(x_0,r)\subset U$.
Finally $S p=p$, so $T^k p=p$ and $p$ is a periodic point of $T$.

Every nonempty open set contains such a point, hence $\Per(T)$ is dense.
\end{proof}

\begin{corollary}
Under the same closed-range hypothesis, every hypercyclic Banach-space
operator with dense small periodic sets is chaotic.
\end{corollary}

\begin{proof}
By the theorem, the periodic points are dense. Adding hypercyclicity is exactly
Devaney chaos in the linear-dynamics convention used by the source paper.
\end{proof}

\section*{Verification And Novelty Notes}

This packet is a partial result, not a full solution of Question 4.2. It covers
operators for which all periodic defect operators $I-T^n$ have closed range.
The proof gives no conclusion when the relevant $I-T^k$ has nonclosed range;
indeed, the mechanism then produces only approximate fixed points for $T^k$.

Local duplicate checks were run in the lightweight indexes for arXiv:2003.07161,
``dense small periodic sets'', ``closed range'', ``approximate periodic'', and
``Cesaro''. No existing packet for this closed-range subcase was found. The
source paper itself proves the weak-star compact and weighted-shift cases, but
does not state this closed-range quotient argument.

The attempted Lipschitz-free counterexample route was checked against
Abbar--Coine--Petitjean, \emph{On the dynamics of Lipschitz operators},
arXiv:2011.10800. That route is not promoted here: their examples show that
linearization may create dense periodic vectors even when the base map has
non-dense periodic points.

\section*{References}

\begin{enumerate}
\item R. Cardeccia and S. Muro, \emph{Arithmetic progressions and chaos in
linear dynamics}, arXiv:2003.07161.
\item M. Abbar, Q. Coine, and C. Petitjean, \emph{On the dynamics of Lipschitz
operators}, arXiv:2011.10800.
\end{enumerate}

\end{document}
